% libre-compendium
%
% Copyright (C) 2026 Richard Qin
%
% This document is dual-licensed under either:
% - GNU Free Documentation License, Version 1.3 or any later version
%   published by the Free Software Foundation; with no Invariant Sections,
%   no Front-Cover Texts, and no Back-Cover Texts (see COPYING.DOC); or
% - Creative Commons Attribution-ShareAlike 4.0 International License
%   (see LICENSE-CC-BY-SA-4.0).
\ProvidesFile{ch01-overview.tex}[2026/11/26 v0.1.0 第一章习题]
\section{数据科学概论}

\begin{exercise}[title={数据驱动求解流程},level=basic]
简述探索性数据分析（EDA）在数据科学问题求解流程中的作用。
\begin{answer}
帮助理解数据分布、发现缺失与异常值、识别变量关系并为特征工程与建模假设提供依据。
\end{answer}
\begin{solution}
在正式建模前，通过描述性统计与可视化检查数据质量（缺失值、重复值、异常值）及特征分布形态，能够避免盲目套用模型，并为数据清洗、特征变换与模型选择提供直接依据。
\end{solution}
\end{exercise}
